\section{Sprintalapú működés}

A sprintalapú működés célja, hogy a fejlesztés ne egyetlen nagy, nehezen követhető feladatként jelenjen meg, hanem kisebb, jól ellenőrizhető iterációkra bontva haladjon. A WatchDB esetében ez azért különösen hasznos, mert az MVP és a későbbi verziók már a PRD-ben is elkülönülnek egymástól. Így a sprintek nem önkényesen kijelölt időszakok, hanem a terméktervezés során meghatározott prioritások gyakorlati megvalósításai.

A folyamat első lépése egy előkészítő sprint, vagyis a Sprint 0 lehet. Ebben nem feltétlenül készül még látványos felhasználói funkció, viszont létrejönnek azok az alapok, amelyekre a későbbi fejlesztés épülhet: a projektstruktúra, az adatbázis kezdeti terve, az autentikációs irány, a CI/CD alapjai, valamint a kezdeti backlog. Ez a szakasz azért fontos, mert csökkenti annak kockázatát, hogy a fejlesztés később szervezetlen technikai döntésekre épüljön.

Az MVP-hez kapcsolódó sprintek már konkrét felhasználói értéket adó funkciókra koncentrálnak. Ilyen lehet például a regisztráció és bejelentkezés elkészítése, a filmkeresés és filmrészletek megjelenítése, a watchlist kezelése, a megnézett filmek naplózása, az értékelés, valamint a saját profiloldal kialakítása. Ezek a funkciók egymásra épülnek, ezért a sprinttervezésnél fontos, hogy a technikai függőségek és a felhasználói folyamatok sorrendje is figyelembe legyen véve.

A későbbi sprintek a termék közösségi és haladó funkcióira épülhetnek. A PRD alapján a v1.1 verzióban jelenhetnek meg a publikus profilok, a követési rendszer, az aktivitási feed és a reakciók, míg a v1.2 már mélyebb funkciókat, például egyéni listákat, statisztikákat, kommenteket vagy ajánlásokat tartalmazhat. Ezek nem tartoznak az MVP alapvető működéséhez, de a sprintalapú tervezés lehetővé teszi, hogy később kontrollált módon kerüljenek be a fejlesztési folyamatba.

Minden sprint végén fontos valamilyen ellenőrizhető eredmény. Ez lehet egy működő feature, egy javított folyamat, egy sikeres build vagy egy tesztelhető preview deployment. A sprint lezárása után a visszajelzések alapján frissíthető a backlog, módosíthatók a prioritások, és pontosítható a következő sprint célja. Így a fejlesztési folyamat nem merev tervként, hanem folyamatosan finomítható munkamenetként működik.
